The Euclidean norm is continuous, so its bounded image makes $X$ bounded. If $a\in\overline X\setminus X$, the continuous function $x\mapsto1/\|x-a\|$ on $X$ is unbounded. Thus $X$ is closed, and Heine--Borel gives compactness.

The same assertion holds in every metric space. First suppose $(x_n)$ is Cauchy without a limit in $X$. The function $g(x)=\lim_n d(x,x_n)$ exists since the real sequence is Cauchy by the reverse triangle inequality. It is $1$-Lipschitz and strictly positive everywhere, since a zero would be a limit of $(x_n)$. The Cauchy property gives $g(x_n)\to0$, so $1/g$ is continuous and unbounded. Hence the bounded-image assumption forces completeness.

If $X$ is not totally bounded, some $\epsilon>0$ permits a sequence $(z_n)$ with pairwise distances at least $\epsilon$, chosen recursively outside finitely many $\epsilon$-balls. Put $\varphi_n(x)=\max(0,1-4d(x,z_n)/\epsilon)$. A ball of radius $\epsilon/8$ about any point meets at most one nonzero set of these functions, since otherwise the two centres would be less than $3\epsilon/4$ apart. Therefore $x\mapsto\sum_{n\ge1} n\varphi_n(x)$ is locally a finite sum of continuous functions, yet takes value $n$ at $z_n$. This contradiction forces total boundedness. A complete totally bounded metric space is compact. The empty space also satisfies the assertion.
